\documentclass[11pt]{article}
\usepackage[T1]{fontenc}
\usepackage[utf8]{inputenc}
\usepackage{lmodern,amsmath,amssymb,amsthm,mathtools,booktabs,longtable}
\usepackage[margin=1in]{geometry}
\usepackage{microtype}
\usepackage[hidelinks]{hyperref}
\newtheorem{theorem}{Theorem}[section]
\newtheorem{proposition}[theorem]{Proposition}
\newtheorem{lemma}[theorem]{Lemma}
\theoremstyle{definition}\newtheorem{definition}[theorem]{Definition}
\theoremstyle{remark}\newtheorem{remark}[theorem]{Remark}
\newcommand{\C}{\mathbb C}\newcommand{\Z}{\mathbb Z}\newcommand{\R}{\mathbb R}
\DeclareMathOperator{\tr}{tr}\DeclareMathOperator{\Rea}{Re}
\newcommand{\status}{\textbf{Research draft. GRH is not proved here.}}
\date{November 2, 2026}

\title{Obstructions and a Corrected Spectral Criterion for GRH}
\author{GRH Research Project\\\small Working manuscript; authorship to be confirmed}
\begin{document}\maketitle\begin{center}\status\end{center}
\begin{abstract}
We audit proposed surface-order and Hilbert--P\'olya arguments for GRH.
We prove obstructions to the proposed prime/geodesic identification,
integral cyclic induction, absolute prime-tail estimates and determinant
preservation, and retain a conditional spectral closure criterion.
We construct convergent arithmetic approximants with explicit errors and
analyze coefficient-defined Jensen metrics. From the explicit scalar
local factors of a separate exponent argument, we extend an Euler-correction
region and its principal-residue interface, and isolate a shifted infinite-order angular Hecke factor.
Under the stated classical analytic inputs, its zeros obstruct extending
the unchanged correction's nonvanishing criterion below $7/12$.
We also reconstruct a frozen-capacity row count conditionally on the
source moment and witness estimates, and extend the low interface under
explicit reflected-energy and Gram hypotheses. The nonprincipal central
contours also extend under their stated arithmetic and buffered-bound
contracts. We verify the finite-field Gauss and zero-preserving correlation
input, with prime-power and reciprocity dependencies distinguished.
We also prove an annular periodic cancellation lemma and reconstruct the
Gram summation under explicit arithmetic representation and period inputs.
A classical conductor argument supplies an explicit common arithmetic period,
distinguished from the smaller quadratic-phase modulus.
We verify the local reflected coefficient separation and its common-profile
normalization, retaining the global theta and hybrid-sieve dependencies.
A pinned Lean 4 companion verifies
selected arithmetic, complex-analytic,
matrix and infinite-product components, with its remaining dependencies
specified. These results do not establish GRH; the full arithmetic
spectral bridge and the required mean-square and contour estimates
remain unresolved.
\end{abstract}
\section{Scope and conventions}
For a primitive Dirichlet character $\chi$ of conductor $q$, write
\[
 L(s,\chi)=\sum_{n\ge1}\frac{\chi(n)}{n^s},\qquad
 \Lambda(s,\chi)=\left(\frac q\pi\right)^{(s+a)/2}
 \Gamma\left(\frac{s+a}{2}\right)L(s,\chi),\quad
 a=\frac{1-\chi(-1)}2.
\]
The series and Euler product initially apply for $\Rea s>1$.
Our target is Dirichlet GRH: every zero in $0<\Rea s<1$ has real part $1/2$.
For $q=1$, use the entire function
$\xi(s)=\tfrac12s(s-1)\Lambda(s,1)$ when an entire target is required.
Finite-order Hecke characters and general Artin representations require separate constructions;
a theorem for one family is not silently a theorem for another.

The supplied surface-order manuscripts \cite{surface,grhsource} propose
arithmetic transfer operators and self-adjoint approximants. We isolate
valid closure logic and prove obstructions to particular claimed constructions.
These obstructions refute intermediate statements, not GRH.
The separate \texttt{math} repository \cite{mathsource} includes claimed zero-free
regions strictly weaker than GRH. Such a region cannot replace the missing critical-line argument.

\section{A corrected conditional closure theorem}
\begin{theorem}[Spectral approximation criterion]\label{thm:closure}
Let $U\subset\C$ be connected and open, let $f$ be holomorphic on $U$ and
not identically zero, and let
\[
 F_N(s)=e^{a_N+b_Ns}\det\bigl(sI-(\tfrac12 I+iH_N)\bigr),
 \qquad H_N=H_N^*\in M_{d_N}(\C).
\]
If $F_N\to f$ locally uniformly on $U$, then every zero of $f$ in $U$
has real part $1/2$.
\end{theorem}
\begin{proof}
The eigenvalues $\lambda_{j,N}$ of $H_N$ are real. Thus the zeros of $F_N$
are $1/2+i\lambda_{j,N}$, with algebraic multiplicities; the exponential
factor never vanishes. If $f(z)=0$ off that line, choose a closed disk
$D\subset U$ about $z$ disjoint from the line, with $f$ nonzero on its
boundary. Such a disk exists because the zeros of a nonzero holomorphic
function on a connected domain are isolated. Let
$m=\min_{\partial D}|f|>0$. Local uniform convergence gives
$|F_N-f|<m$ on $\partial D$ for sufficiently large $N$.
Rouch\'e's theorem forces $F_N$ and $f$ to have the same positive number
of zeros inside $D$, contradicting the zero location of $F_N$.
\end{proof}

\begin{remark}
This is a conditional criterion, not a construction. Self-adjointness alone
establishes only the zero location of the approximants. The difficult arithmetic
statement is convergence to the particular completed $L$-function.
Pointwise convergence, weak operator convergence, and agreement of finitely
many zeros do not supply the displayed hypothesis.
\end{remark}

\begin{proposition}[A usable continuation interface]\label{prop:normal}
Suppose $U$ is a connected domain, the entire functions $F_N$ are locally
uniformly bounded on $U$, and they converge pointwise on a nonempty open
subset $V\subset U$ to a holomorphic function $f$ defined on $U$.
Then $F_N\to f$ locally uniformly on $U$.
\end{proposition}
\begin{proof}
Montel's theorem makes the family normal. Every subsequence has a further
subsequence converging locally uniformly to a holomorphic function $g$.
On $V$, pointwise convergence gives $g=f$, so the identity theorem gives
$g=f$ on $U$. If the original sequence failed to converge locally uniformly,
a compact set, positive error, and violating subsequence would contradict
this subsequence property.
\end{proof}
A proposed analytic continuation proof must actually provide the uniform
bounds and a genuine overlap $V$. The disjoint domains $\Rea s>1$ and
$0<\Rea s<1$ are not an overlap.

\section{Arithmetic obstructions to the claimed completion}
\begin{theorem}[No integral prime-power geodesic norms]\label{thm:norm}
Let $\gamma\in\mathrm{SL}_2(\Z)$ be hyperbolic. Its expanding eigenvalue
in absolute value is $u>1$, and its hyperbolic norm is $N(\gamma)=u^2$.
Then $N(\gamma)$ is not an integer greater than one. In particular no
closed geodesic arising from such a matrix has length $k\log p$ for a
prime $p$ and integer $k\ge1$.
\end{theorem}
\begin{proof}
Put $t=|\tr\gamma|\in\Z$, so $t=u+u^{-1}$.
If $u^2=n\in\Z$, $n>1$, squaring gives
\[
 t^2=n+2+\frac1n.
\]
The left-hand side and $n+2$ are integers, so $1/n$ is an integer,
contradicting $0<1/n<1$. Since length is $\log N(\gamma)$,
length $k\log p$ would imply $N(\gamma)=p^k$, already excluded.
\end{proof}
The statement applies to every subgroup of $\mathrm{SL}_2(\Z)$, hence to
congruence surfaces. The prime geodesic theorem is an asymptotic counting
analogy; it cannot provide the termwise identification claimed in the
completion section of \cite{surface}. Hecke correspondences of determinant
$p$ must not be identified with hyperbolic deck transformations of determinant one.

\begin{theorem}[Failure of integral cyclic induction]\label{thm:cyclic}
The trivial character of $G=C_2\times C_2$ is not an integer linear
combination of characters induced from one-dimensional characters of
cyclic subgroups of $G$.
\end{theorem}
\begin{proof}
Every cyclic subgroup has order one or two. The degree of any such
induced character is its index, namely four or two. Evaluation at the
identity sends every integer linear combination to an even integer,
whereas the trivial character has degree one.
\end{proof}
This distinguishes rational Artin induction from integral Brauer induction.
It invalidates the stated cyclic integral reduction, while leaving appropriate
forms of these classical induction theorems available for a revised argument.

\begin{theorem}[Prime-tail obstruction]\label{thm:tail}
For every real $\sigma\le1$, the positive series
$\sum_p p^{-\sigma}$ diverges. Thus
$\sum_{p>N}p^{-\sigma}$ is not a vanishing absolutely convergent tail
in the critical strip.
\end{theorem}
\begin{proof}
It suffices to show $\sum_p1/p=\infty$, since $p^{-\sigma}\ge p^{-1}$.
Suppose the reciprocal-prime series converged. For $0\le x\le1/2$,
$-\log(1-x)\le2x$, for example by integrating
$(1-t)^{-1}\le2$ from zero to $x$. Hence the finite Euler products
$\prod_{p\le X}(1-1/p)^{-1}$ would be bounded uniformly in $X$.
Expanding each as a finite product of absolutely convergent geometric
series gives the sum of $1/n$ over all positive integers with prime
factors at most $X$. In particular it dominates
$\sum_{n\le\lfloor X\rfloor}1/n$, which is unbounded. Contradiction.
\end{proof}
Any repair below $\Rea s=1$ must introduce proved cancellation or a
regularization, with quantitative control; absolute prime-tail bounds cannot suffice.

\section{Determinants and lost arithmetic information}
For $T=\operatorname{diag}(\chi(p)p^{-s})_{p\le N}$ one has exactly
\[
 \det(I-T)=\prod_{p\le N}(1-\chi(p)p^{-s})=L_N(s,\chi)^{-1},
\]
not $L_N$. On an absolute-convergence region,
\[
 \log\det(I-T)=-\sum_{k\ge1}\frac{\tr(T^k)}k.
\]
The reciprocal orientation matters before any analytic continuation.
Moreover Hermitian-part extraction does not preserve this determinant:
\[
 T=\begin{pmatrix}0&1\\-1&0\end{pmatrix},\quad
 \frac{T+T^*}2=0,\quad \det(I-T)=2\ne1=\det(I).
\]
An identity for the transfer determinant therefore does not automatically
identify the centered characteristic polynomial of a symmetrized operator.

\begin{proposition}[A direct obstruction to the supplied diagonal diagnostic]\label{prop:diagonal}
For the $q=1$ model evaluated at $s=1/2$, the algorithm that sets
$\lambda_p=-\arg(p^{-1/2})$ gives $H_N=0$ and
$F_N(s)=(s-1/2)^{d_N}$. These functions cannot converge locally uniformly
to $\Lambda(s,1)$ or $\xi(s)$ on a neighborhood of $1/2$.
\end{proposition}
\begin{proof}
Every $p^{-1/2}$ is positive, so its argument is zero and $F_N(1/2)=0$.
On the other hand the alternating Dirichlet series
$\eta(1/2)=\sum_{n\ge1}(-1)^{n-1}n^{-1/2}$ is positive: grouping into
adjacent pairs gives positive terms and its sum is at least
$1-1/\sqrt2>0$. The continuation identity
$\eta(s)=(1-2^{1-s})\zeta(s)$ gives $\zeta(1/2)<0$.
The gamma factor at $1/4$ is positive and finite, so
$\Lambda(1/2,1)\ne0$, and its pole-removal factor at $1/2$ is nonzero.
Even pointwise convergence at that point is impossible.
\end{proof}

\begin{proposition}[Holomorphic Hermitian rigidity]\label{prop:hermitian-rigidity}
Let $U\subset\C$ be connected and open, with a fixed inner product on
$\C^d$. If $H:U\to M_d(\C)$ has holomorphic entries and
$H(s)^*=H(s)$ for every complex $s\in U$, then $H$ is constant.
\end{proposition}
\begin{proof}
In fixed orthonormal coordinates, for each vector $v$, the holomorphic
scalar $v^*H(s)v$ is real-valued,
so it is constant by the open mapping theorem. Polarization recovers all
matrix entries from these quadratic forms, hence every entry is constant.
Equivalently, $H_{ij}+H_{ji}$ and $i(H_{ij}-H_{ji})$ are real-valued
holomorphic functions whose constancy determines $H_{ij}$.
\end{proof}
Thus a nonconstant arithmetic transfer $T(s)$ cannot retain entrywise
holomorphy after pointwise Hermitian-part extraction: a holomorphic
$(T+T^*)/2$ would be constant, and the remaining holomorphic skew-Hermitian
part would be constant as well. This excludes that particular shortcut,
not fixed self-adjoint operators. The spectral criterion uses $H_N$
independent of the spectral variable $s$. Requiring self-adjointness only
on a real parameter line is also consistent with a nonconstant holomorphic
extension, but needs its own zero-location theorem.

\begin{proposition}[A positive generalized spectral pencil]\label{prop:positive-pencil}
Let $A=A^*$ and $B=B^*>0$ be fixed $d\times d$ matrices.
Every zero of $\det(A-zB)$ is real. Consequently every zero of
$\det(A+i(s-1/2)B)$ lies on $\Rea s=1/2$.
\end{proposition}
\begin{proof}
Singularity supplies $v\ne0$ with $Av=zBv$. Hence
$z=(v^*Av)/(v^*Bv)$, with real numerator and strictly positive real
denominator. Thus $z$ is real. Substitute $z=-i(s-1/2)$.
\end{proof}
This supplies a valid finite spectral mechanism, but not its arithmetic
construction. To apply it one must construct $A_N,B_N$, certify
$B_N>0$, and prove the required completed-$L$ determinant approximation
independently. Hermitian coefficients alone are insufficient: for
\[
 A=\begin{pmatrix}0&\varepsilon\\\varepsilon&0\end{pmatrix},
 \quad B=\begin{pmatrix}1&0\\0&-1\end{pmatrix},\qquad
 \det(A-zB)=-(z^2+\varepsilon^2),
\]
$0<\varepsilon<1/2$ gives transported off-line zeros
$s=1/2\pm\varepsilon$ inside the critical strip.

\section{A necessary change of approximation class}
\begin{theorem}[Zero-free Euler truncations cannot approximate a zero]\label{thm:euler-nogo}
Let $U$ be a connected open subset of $\{\Rea s>0\}$, let $g_N$
be holomorphic and nowhere zero on $U$, and let $P_N$ be arbitrary
finite sets of rational primes. Put
\[
 D_N(s)=\prod_{p\in P_N}(1-\chi(p)e^{-s\log p}).
\]
Each of $g_ND_N$ and $g_N/D_N$ is holomorphic and nowhere zero on $U$.
If either sequence converges locally uniformly to a holomorphic $f$
which is not identically zero, then $f$ has no zero in $U$.
\end{theorem}
\begin{proof}
For $\Rea s>0$, $|\chi(p)e^{-s\log p}|\le p^{-\Rea s}<1$,
so every local factor is nonzero. Hence the finite product, its reciprocal,
and their normalized versions have the asserted property.
If $f$ had a zero, choose a closed disk inside $U$ with a positive
number of zeros of $f$ in its interior and none on its boundary.
Local uniform convergence and Rouch\'e's theorem would force the
zero-free approximants to have that positive number of zeros, a contradiction.
\end{proof}
More quantitatively, on any such disk $D$, every zero-free holomorphic
approximant $G$ satisfies
\[
 \max_{\partial D}|G-f|\ \ge\ \min_{\partial D}|f|>0.
\]
Otherwise Rouch\'e yields the same contradiction. Thus improved estimates
cannot repair this approximation class near a nontrivial zero. The
construction must change so that finite approximants can have zeros in
the strip; a nonvanishing gamma or exponential normalization cannot do this.

\section{A convergent arithmetic replacement}
For a nonprincipal character modulo $q$, put
$A(x)=\sum_{n\le x}\chi(n)$ and $B=\sup_x|A(x)|\le q$.
Complete periods sum to zero, so $B$ is finite. Partial summation gives,
for $\sigma=\Rea s>0$,
\[
 L(s,\chi)-\sum_{n\le N}\chi(n)n^{-s}
 =-A(N)N^{-s}+s\int_N^\infty A(x)x^{-s-1}\,dx,
\]
and therefore
\[
 \left|L(s,\chi)-\sum_{n\le N}\chi(n)n^{-s}\right|
 \le B\left(1+\frac{|s|}{\sigma}\right)N^{-\sigma}.
\]
The endpoint tends to zero and the integral is absolutely convergent.
The limit agrees with the Euler definition on $\Rea s>1$, hence with
its analytic continuation. Multiplication by the correct gamma/conductor
factor gives locally uniform approximation on $\Rea s>0$ for either parity.

For zeta, reproduce the pole before removing it. Set
\[
 Z_N(s)=\sum_{n\le N}n^{-s}+\frac{N^{1-s}}{s-1}.
\]
Euler summation gives, for $\sigma>0$ and $s\ne1$,
\[
 \zeta(s)-Z_N(s)=-s\int_N^\infty\{x\}x^{-s-1}\,dx,
 \qquad |\zeta(s)-Z_N(s)|\le\frac{|s|}{\sigma}N^{-\sigma}.
\]
The integral continues holomorphically to $\sigma>0$. Consequently
\[
 \xi_N(s)=\frac{s}{2}\pi^{-s/2}\Gamma(s/2)
 \left((s-1)\sum_{n\le N}n^{-s}+N^{1-s}\right)
\]
is holomorphic on that half-plane, including $s=1$, and converges
locally uniformly to $\xi$; in fact $\xi_N(1)=1/2$ for every $N$.
These constructions use known arithmetic coefficients, not unknown zeros.
They avoid Theorem~\ref{thm:euler-nogo} by allowing additive cancellation
and unrestricted finite-level zeros. They establish convergence, but
provide no self-adjoint realization or critical-line theorem for those zeros.

\section{A certified test of the arithmetic pencil candidate}
Let $\chi$ be the primitive nonprincipal character modulo $3$. Its fifth
Dirichlet partial sum is the entire function
\[
 f(s)=S_5(s,\chi)=1-2^{-s}+4^{-s}-5^{-s}.
\]
\begin{theorem}[Certified zero of a finite partial sum]\label{thm:finite-offline}
Put
\[c=(337261+11027504i)/10^6,\qquad r=10^{-4}.\]
The disk $|s-c|<r$ contains exactly one zero of $f$, counted with
multiplicity, and lies inside $0<\Rea s<1/2$.
\end{theorem}
\begin{proof}
The exact rational procedure detailed in the extended edition establishes
\[
 |f(c)|<10^{-6},\qquad |f'(c)|>1,\qquad
 \sup_{|s-c|\le r}|f''(s)|\le9.
\]
It uses $60$ terms of the logarithm's positive atanh series and terms
$0$ through $120$ of the exponential Taylor series, with explicit
remainders; no floating-point value is used in these inequalities.
On the boundary, Taylor's formula gives
\[
 |f(s)-f'(c)(s-c)|<10^{-6}+\tfrac92r^2<r<|f'(c)(s-c)|.
\]
Rouch\'e gives exactly one zero, which is therefore simple.
The real part throughout the closed disk belongs to
$[337161,337361]/10^6\subset(0,1/2)$.
\end{proof}
The completed fifth partial sum has the same zero because its
conductor/gamma factor is holomorphic and nonzero on this disk.
It cannot equal a nowhere-zero holomorphic factor times the determinant of a fixed
positive Hermitian pencil there. This rejects that exact finite-level
identification. It does not reject convergence of the partial sums,
modified approximants, later truncations, subsequences, or GRH itself:
$f$ is a finite partial sum, not the full $L$-function.

\section{Jensen approximants and a coefficient-defined metric}
This alternative currently concerns the Riemann target only. Define
\[
 \gamma_j=\frac{j!\,\xi^{(2j)}(1/2)}{(2j)!},\qquad
 F(w)=\sum_{j\ge0}\frac{\gamma_j}{j!}w^j,\qquad
 \Xi(t)=\xi(1/2+it)=F(-t^2).
\]
The classical positive Fourier kernel of $\Xi$ gives $\gamma_j>0$~\cite{romik};
evenness and entirety of $\xi$ make $F$ entire. Set, for $d\ge1$,
\[
 J_d(x)=\sum_{j=0}^d\binom dj\gamma_jx^j,\qquad
 Q_d(w)=J_d(w/d),\qquad P_d(t)=Q_d(-t^2).
\]
\begin{proposition}[Explicit convergence without a zero-location premise]
For $R,T\ge0$,
\[
 \sup_{|w|\le R}|Q_d(w)-F(w)|\le\frac{R^2F''(R)}{2d},\qquad
 \sup_{|t|\le T}|P_d(t)-\Xi(t)|\le\frac{T^4F''(T^2)}{2d}.
\]
\end{proposition}
\begin{proof}
The coefficient multiplier $w_{d,j}=j!\binom dj/d^j$ for $j\le d$,
and $0$ for $j>d$, satisfies
$0\le w_{d,j}\le1$ and $1-w_{d,j}\le j(j-1)/(2d)$.
For $j\le d$ use $1-\prod_{k<j}(1-k/d)\le\sum_{k<j}k/d$;
for $j>d$ use $j(j-1)\ge2d$. Summing the absolutely convergent
series with nonnegative coefficients proves the bounds.
\end{proof}

Let $C_d$ be the multiplication-by-$x$ companion of the monic
$J_d/\gamma_d$, and $H_d=(m_{i+j})_{0\le i,j<d}$, where
$m_k=\tr(C_d^k)$ and $m_0=d$. Newton identities compute $H_d$
from the coefficients, and $H_dC_d=C_d^T H_d$.
The Hermite--Sylvester criterion gives $H_d\ge0$ exactly when
$J_d$ is real-rooted~\cite{hermitesylvester}; then its positive coefficients
force all roots negative. If $H_d>0$, the extended edition constructs
an explicit real symmetric $K_d$ with
\[
 P_d(t)=\frac{(-1)^d\gamma_d}{d^d}\det(tI-K_d).
\]
Strict positivity is needed for that construction; semidefinite metrics
can be singular. The classical Jensen criterion makes $H_d\ge0$ for
every $d$ equivalent to zeta RH~\cite{jensen}. This identifies the missing
positivity; it does not prove it. Fixed-degree eventual hyperbolicity
as the shift grows does not settle this unshifted growing-degree family.

\begin{proposition}[Positive smooth kernels do not supply the metric]
Positivity of an even smooth Fourier density, including Gaussian decay,
does not imply positivity of its Jensen Newton-sum metric.
\end{proposition}
\begin{proof}
Convolve $\frac9{10}\delta_0+\frac1{20}(\delta_{-1}+\delta_1)$ with a
centered Gaussian of variance $v>0$. Its Fourier transform is
$e^{-vt^2/2}(9/10+\cos t/10)$ and its density is strictly positive,
even and Schwartz. Writing its associated $F(w)=\sum\gamma_jw^j/j!$
gives $\gamma_0=1$, $\gamma_1=(v+1/10)/2$ and
$\gamma_2=(3v^2+3v/5+1/10)/12$. Thus
$\gamma_1^2-\gamma_0\gamma_2=-7/1200<0$; the degree-two
Newton-sum metric has negative determinant. These are model coefficients,
not the arithmetic $\xi$ coefficients.
\end{proof}

\section{Consequences for a revised research article}
Theorems~\ref{thm:norm}, \ref{thm:cyclic}, and \ref{thm:tail}
require replacements for the associated proof steps, not additional exposition.
A credible next theorem would construct operators independently of unknown
zeros and establish their exact arithmetic determinant, including ramified,
archimedean, and conductor terms. It must then prove locally uniform
convergence or the local boundedness in Proposition~\ref{prop:normal}.
Only then does Theorem~\ref{thm:closure} prove the desired zero location.

Weil/Li positivity, adelic trace methods, and canonical-product approximations
are potential alternative interfaces. They move the research obligation to a
specific positivity, trace, or approximation theorem; they do not discharge it
by analogy. Lagarias \cite{lagarias} and Connes \cite{connes} provide relevant
primary-source starting points. No novelty or complete proof is claimed for
these classical conditional principles.

\begin{thebibliography}{9}
\bibitem{surface} A. Schreiber, \emph{Emergent Cyclic Time from Surface Orders and a Hilbert--P\'olya Program for the Generalized Riemann Hypothesis}, repository source,
commit \texttt{212ff730}, \url{https://github.com/amelie-iska/Surface-Orders-Hilbert-Polya}.
\bibitem{grhsource} A. Schreiber, \emph{Surface-Order Determinants, Monomial Artin L-Functions, and the General Hilbert--P\'olya Extension Problem}, repository PDF,
commit \texttt{0c77122d}, \url{https://github.com/amelie-iska/GRH}.
\bibitem{mathsource} \emph{Mathematical manuscripts and supporting proof artifacts}, supplied repository, \url{https://github.com/amelie-iska/math}. See project source manifest for the exact local pin and review scope.
\bibitem{lagarias} J. C. Lagarias, Li coefficients for automorphic $L$-functions,
\emph{Ann. Inst. Fourier} \textbf{57} (2007), 1689--1740,
\url{https://arxiv.org/abs/math/0404394v4}.
\bibitem{connes} A. Connes, Trace formula in noncommutative geometry and the zeros
of the Riemann zeta function, \emph{Selecta Math.} \textbf{5} (1999), 29--106,
\url{https://arxiv.org/abs/math/9811068}.
\bibitem{mathlib} The mathlib community, \emph{Riemann zeta function documentation},
\url{https://leanprover-community.github.io/mathlib4_docs/Mathlib/NumberTheory/LSeries/RiemannZeta.html}.
\end{thebibliography}
\end{document}
